\section*{Розділ VI. Виборчі процедури}
\addcontentsline{toc}{section}{Розділ VI. Виборчі процедури}

\subsection*{6.1. Етапи виборчого процесу}
\addcontentsline{toc}{subsection}{6.1. Етапи виборчого процесу}
    6.1.1. Виборчий процес включає такі обов'язкові етапи:

        \begin{enumerate}[label=\alph*)]
            \item Призначення виборів;
            \item Висування кандидатів;
            \item Реєстрація кандидатів;
            \item Передвиборча агітація;
            \item Голосування;
            \item Підведення підсумків та оголошення результатів.
        \end{enumerate}

    6.1.2. Кожен етап має визначені строки та процедури, які встановлюються ЦВКс відповідно до цього Положення.

    6.1.3. Перехід до наступного етапу можливий лише після завершення попереднього етапу у встановлені строки.

\subsection*{6.2. Призначення виборів}
\addcontentsline{toc}{subsection}{6.2. Призначення виборів}
    6.2.1. Вибори до органів студентського самоврядування призначаються ЦВКс не пізніше ніж за 45 календарних днів до дати закінчення повноважень відповідного органу або посадової особи.

    6.2.2. Позачергові вибори призначаються ЦВКс протягом 7 календарних днів з дня настання підстав для їх проведення.

    6.2.3. Рішення про призначення виборів повинно містити:

        \begin{enumerate}[label=\alph*)]
            \item Дату та час проведення голосування;
            \item Посади або склад органу, що обираються;
            \item Календарний план виборчого процесу з зазначенням строків кожного етапу;
            \item Перелік виборчих дільниць та їх межі;
            \item Порядок висування та реєстрації кандидатів;
            \item Способи голосування (паперове, електронне або змішане).
        \end{enumerate}

    6.2.4. Рішення про призначення виборів оприлюднюється на офіційних інформаційних ресурсах органів студентського самоврядування не пізніше ніж наступного дня після його прийняття.

\subsection*{6.3. Строки виборчого процесу}
\addcontentsline{toc}{subsection}{6.3. Строки виборчого процесу}
    6.3.1. Загальна тривалість виборчого процесу становить 30 календарних днів з дня призначення виборів до дня оголошення результатів.

    6.3.2. Розподіл строків за етапами:

        \begin{enumerate}[label=\alph*)]
            \item Висування кандидатів -- 7 календарних днів;
            \item Реєстрація кандидатів -- 5 календарних днів;
            \item Передвиборча агітація -- 14 календарних днів;
            \item Голосування -- 1-3 календарні дні (залежно від способу голосування);
            \item Підведення підсумків -- 3 календарні дні.
        \end{enumerate}

    6.3.3. У виняткових випадках ЦВКс може скоротити строки виборчого процесу, але загальна тривалість не може бути менше 21 календарного дня.

    6.3.4. Строки виборчого процесу для позачергових виборів можуть бути скорочені до 14 календарних днів за рішенням ЦВКс.

\subsection*{6.4. Висування та реєстрація кандидатів}
\addcontentsline{toc}{subsection}{6.4. Висування та реєстрація кандидатів}
    6.4.1. Право висування кандидатів мають:

        \begin{enumerate}[label=\alph*)]
            \item Органи студентського самоврядування;
            \item Студентські організації, зареєстровані в Університеті;
            \item Ініціативні групи студентів у складі не менше 15 осіб для виборів Голови СР КАІ, 10 осіб для виборів Голови СРФ/СРІ;
            \item Студенти можуть висувати свою кандидатуру самостійно (самовисування).
        \end{enumerate}

    6.4.2. Для реєстрації кандидат подає до ЦВКс:

        \begin{enumerate}[label=\alph*)]
            \item Заяву про згоду балотуватися;
            \item Довідку про статус студента Університету;
            \item Програму кандидата (не більше 2 сторінок);
            \item Документ про висування (крім самовисування).
        \end{enumerate}

    6.4.3. ЦВКс приймає рішення про реєстрацію або відмову в реєстрації кандидата протягом 3 робочих днів з дня подання документів.

    6.4.4. Підставами для відмови в реєстрації є:

        \begin{enumerate}[label=\alph*)]
            \item Невідповідність кандидата встановленим вимогам;
            \item Подання неповного пакету документів;
            \item Подання недостовірних відомостей;
            \item Порушення процедури висування.
        \end{enumerate}

\subsection*{6.5. Передвиборча агітація}
\addcontentsline{toc}{subsection}{6.5. Передвиборча агітація}
    6.5.1. Передвиборча агітація -- це діяльність, спрямована на спонукання виборців голосувати за або проти певного кандидата.

    6.5.2. Передвиборча агітація дозволяється з дня реєстрації кандидатів і припиняється о 24:00 дня, що передує дню голосування.

    6.5.3. ЦВКс забезпечує рівні можливості для агітаційної діяльності всіх зареєстрованих кандидатів, включаючи:

        \begin{enumerate}[label=\alph*)]
            \item Рівний доступ до офіційних інформаційних ресурсів;
            \item Рівні можливості проведення агітаційних заходів;
            \item Рівний час для представлення програм на офіційних зустрічах з виборцями.
        \end{enumerate}

    6.5.4. Забороняється:

        \begin{enumerate}[label=\alph*)]
            \item Підкуп виборців або кандидатів;
            \item Використання службового становища для агітації;
            \item Поширення недостовірної інформації про кандидатів;
            \item Агітація в день голосування та на виборчих дільницях.
        \end{enumerate}

\subsection*{6.6. Способи та порядок голосування}
\addcontentsline{toc}{subsection}{6.6. Способи та порядок голосування}
    6.6.1. Голосування може проводитися у таких формах:

        \begin{enumerate}[label=\alph*)]
            \item Паперове голосування -- з використанням паперових бюлетенів на виборчих дільницях;
            \item Електронне голосування -- з використанням електронних систем голосування;
            \item Змішане голосування -- з можливістю вибору способу голосування виборцем.
        \end{enumerate}

    6.6.2. Спосіб голосування визначається ЦВКс у рішенні про призначення виборів з урахуванням технічних можливостей та особливостей конкретних виборів.

    6.6.3. При електронному голосуванні забезпечується:

        \begin{enumerate}[label=\alph*)]
            \item Ідентифікація виборців;
            \item Таємність голосування;
            \item Неможливість повторного голосування;
            \item Захист від несанкціонованого доступу;
            \item Можливість перевірки результатів.
        \end{enumerate}

    6.6.4. Детальні технічні вимоги до електронних систем голосування (алгоритми шифрування, протоколи передачі даних, резервні процедури при технічних збоях) визначаються технічними регламентами ЦВКс з дотриманням загальних принципів, встановлених цим Положенням.

    6.6.5. Голосування проводиться у робочі дні з 9:00 до 18:00, у вихідні дні -- з 10:00 до 16:00.

    6.6.6. Виборець голосує особисто, голосування за інших осіб не допускається.

\subsection*{6.7. Спостерігачі за виборчим процесом}
\addcontentsline{toc}{subsection}{6.7. Спостерігачі за виборчим процесом}
    6.7.1. Право бути спостерігачами на виборах мають:

        \begin{enumerate}[label=\alph*)]
            \item Представники кандидатів (по одному від кожного кандидата);
            \item Представники органів студентського самоврядування;
            \item Представники студентських організацій, акредитованих ЦВКс;
            \item Студенти Університету за поданням ініціативних груп у складі не менше 10 осіб.
        \end{enumerate}

    6.7.2. Акредитація спостерігачів здійснюється ЦВКс на підставі письмових заяв, поданих не пізніше ніж за 3 дні до початку голосування.

    6.7.3. Спостерігачі мають право:

        \begin{enumerate}[label=\alph*)]
            \item Присутні під час голосування та підрахунку голосів;
            \item Ознайомлюватися з виборчими документами (крім бюлетенів до завершення голосування);
            \item Робити письмові зауваження щодо порушень виборчих процедур;
            \item Отримувати копії протоколів про результати голосування.
        \end{enumerate}

    6.7.4. Спостерігачі не мають права втручатися у роботу ЦВКс та виборчих комісій, агітувати під час голосування або перешкоджати виборчому процесу.

\subsection*{6.8. Підрахунок голосів та встановлення результатів}
\addcontentsline{toc}{subsection}{6.8. Підрахунок голосів та встановлення результатів}
    6.8.1. Підрахунок голосів здійснюється ЦВКс або уповноваженими нею особами одразу після закінчення голосування.

    6.8.2. При паперовому голосуванні підрахунок проводиться публічно за участю представників кандидатів та спостерігачів.

    6.8.3. При електронному голосуванні результати формуються автоматично системою з обов'язковою перевіркою ЦВКс.

    6.8.4. За результатами підрахунку голосів складається підсумковий протокол, який містить:

        \begin{enumerate}[label=\alph*)]
            \item Загальну кількість виборців;
            \item Кількість виборців, які взяли участь у голосуванні;
            \item Кількість голосів, поданих за кожного кандидата;
            \item Кількість недійсних бюлетенів;
            \item Рішення про результати виборів.
        \end{enumerate}

    6.8.5. Обраним вважається кандидат, який отримав найбільшу кількість голосів виборців, які взяли участь у голосуванні. Мінімальна явка виборців для визнання виборів дійсними не встановлюється.

    6.8.6. У разі отримання однакової кількості голосів двома або більше кандидатами, ЦВКс призначає повторне голосування між цими кандидатами.

    6.8.7. У разі відмови обраного кандидата від посади, повноваження автоматично переходять до кандидата, який посів друге місце за кількістю голосів. Якщо такий кандидат також відмовляється або відсутній, ЦВКс призначає повторні вибори.

\subsection*{6.9. Оголошення результатів виборів}
\addcontentsline{toc}{subsection}{6.9. Оголошення результатів виборів}
    6.9.1. Результати виборів оголошуються ЦВКс не пізніше ніж через 24 години після завершення підрахунку голосів.

    6.9.2. Інформація про результати виборів оприлюднюється на офіційних інформаційних ресурсах органів студентського самоврядування та включає:

        \begin{enumerate}[label=\alph*)]
            \item Підсумковий протокол про результати голосування;
            \item Рішення ЦВКс про підсумки виборів;
            \item Інформацію про порядок оскарження результатів.
        \end{enumerate}

    6.9.3. Документи про результати виборів Голови СР КАІ та Голови СР СМ передаються до КСУ для затвердження відповідно до Положення про ОСС.